\section{Lineage and Context}
\label{sec:otel-weaver-exp-005-collector-to-intake-lineage}

\subsection{2016 Language-Model Lesson}
\label{sec:otel-weaver-exp-005-collector-to-intake-lineage-2016}

The 2016 language-model experiment established the foundational negative result that underlies
the entire dissertation: local text imitation does not conserve consequence. A model that
reproduces surface-level linguistic patterns cannot be tested against the real-world outcomes
those patterns were intended to produce. Prediction without coordination is observationally
indistinguishable from prediction with coordination until a process log is examined. EXP-005
is a direct descendant of this lesson. The OTel Collector pipeline is the mechanism that
converts raw telemetry emissions --- which are, like language-model outputs, surface-level
signals --- into process evidence that can be tested for consequence-conservation. Specifically,
the \texttt{pi.intake.ingested\_at} timestamp injected by \texttt{transform/pi} anchors the
event to a clock that is independent of the emitting model, making the temporal ordering of
events auditable by a process mining engine rather than inferred from model internals.

Without the collector boundary, the dissertation cannot distinguish between a process that
emitted spans in the correct causal order and a process that emitted spans in any order ---
the same failure mode the 2016 experiment exposed at the language level.

\subsection{Enterprise Process Gap}
\label{sec:otel-weaver-exp-005-collector-to-intake-lineage-enterprise-gap}

The enterprise process gap identified in the dissertation lineage is the structural absence of
a lawful event log in production systems that nonetheless emit rich observability data. Modern
enterprises running OpenTelemetry-instrumented services produce millions of spans per hour.
Those spans encode timing, causality, and resource attribution --- the raw material of process
intelligence. Yet no standard pipeline converts OTLP spans into OCEL~2.0-structured process
logs. The gap is not a data-availability gap; it is a law-enforcement gap. Spans are emitted
without mandatory binding to a process instance identifier, without BLAKE3 witness sealing,
and without conformance-court-compatible structure.

EXP-005 is the first experiment in the dissertation to directly address this gap at the
infrastructure layer. Rather than proposing new instrumentation obligations for application
developers, it demonstrates that the gap can be closed at the collector boundary using
standard OTTL processor composition. The \texttt{filter/pi} processor enforces
\texttt{process.pi.instance\_id} presence as a namespace gatekeeper, silently dropping spans
that cannot be bound to a process instance. This is the minimal law-enforcement intervention
that converts an enterprise observability pipeline into a process intelligence intake pipeline.

\subsection{Relationship to the Chatman Equation}
\label{sec:otel-weaver-exp-005-collector-to-intake-lineage-chatman-equation}

The Chatman Equation ($A = \mu(O^*)$) states that autonomous action is a function of an
ontological substrate. In the context of EXP-005, the ontological substrate is the process
instance binding enforced by \texttt{process.pi.instance\_id}. A span without an instance
binding cannot be placed within an object-centric event log; it has no substrate. The
\texttt{filter/pi} processor is the operational implementation of the Chatman Equation's
substrate requirement: it enforces that only spans with a lawful ontological binding are
admitted to the intake pipeline. Spans that lack this binding are not rejected with an error;
they are silently dropped, reflecting the doctrine that unbound telemetry is not refusable ---
it is simply not process evidence.

The \texttt{transform/pi} processor operationalises the $\mu$ (manufacturing) operator of the
equation: it normalises raw spans into the canonical form required by the process intelligence
conformance court, lowercasing activity names and injecting provenance metadata.

\subsection{Relationship to Receipts and Replay}
\label{sec:otel-weaver-exp-005-collector-to-intake-lineage-receipts}

The receipt-and-replay architecture of the dissertation requires that every process intelligence
artifact be sealed with a cryptographic receipt that can be used to replay the manufacturing
pipeline and verify the artifact's provenance. EXP-005 is the upstream boundary of this
architecture: it is where raw OTLP spans are converted into the OCEL~2.0 JSON records that
downstream receipt-manufacturing stages can hash and seal. The \texttt{live-check-intake.manifest.yaml}
references a BLAKE3 receipt store at
\texttt{/Users/sac/process-intelligence/receipts/live-check-intake-receipt.json}. This
receipt does not yet exist --- it is the PARTIAL gap documented in the experiment status.
When the receipt is materialised by a live collector run, EXP-005 will be the first
experiment in the OTel Weaver series to produce a receipt-sealed OCEL~2.0 artifact from live
OTLP traffic, completing the upstream end of the receipt chain.

The Span~C test class (admitted by the collector, rejected by the downstream live-check gate)
is specifically designed to verify that the receipt chain correctly records refusals: a span
that passes the collector boundary but fails the conformance court must produce a sealed
refusal record, not a silent drop. This distinguishes the collector law (admission by
attribute presence) from the court law (admission by full schema conformance).
